A fundamental problem when building machine learning systems is to
predict {\em counterfactuals} about model behavior. For example, scaling laws
\citep{kaplan2020scaling,hashimoto2021model,muennighoff2023scaling}
aim to predict the performance of systems trained with more data and
more compute than
is currently available; interpretability techniques
\citep{kim2018interpretability} predict how models behave under counterfactual
inputs.

Analogously, in this work we study {\em predictive data attribution}
(or {\em datamodeling} \citep{ilyas2022datamodels}),
where the goal is to predict how a model would behave if it
had been trained on a different dataset.
This well-studied problem encompasses, e.g., estimating the effect
(on the resulting trained model's predictions) of modifying a training example
\citep{koh2017understanding}, removing a group of
training examples \citep{koh2019accuracy,bae2022if,park2023trak}, or
adding entire training data sources~\citep{ley2024generalized}.

Predictive data attribution in large-scale settings is challenging: it
requires simulating training a model on a different
dataset without actually training
\citep{guu2023simfluence,ilyas2024data}.
In ``classical'' settings---when learning corresponds to minimizing a
convex loss---statistical tools
like the influence function~\citep{hampel1947influence}
allow us to accurately and
efficiently estimate how different training data choices change trained model
predictions~\citep{rad2018scalable,koh2019accuracy,giordano2019swiss}.
However, in the non-convex settings that are ubiquitous in natural
domains like language/vision, current methods
are less effective. Indeed, the best existing methods produce
estimates that typically (a)
only {\em moderately correlate} with the ground
truth~\citep{basu2021influence,bae2022if,park2023trak} and (b)
incur large absolute error~\citep{bae2022if}.


\subsection{Contributions}
In this work, we make progress towards solving the
well-studied problem of estimating how a model's
predictions would change under different
(counterfactual) deletions of training data
\citep{koh2017understanding,bae2022if,schioppa2022scaling,park2023trak,bae2024training}.
We make two main contributions:

\begin{figure}[ht]
  \centering
  \newcommand{\dataset}{cifar}
  \newcommand{\exampleid}{0}
  \newcommand{\scatteropacity}{0.5}
  \newcommand{\setting}{ResNet-9 on CIFAR-10}
  \newcommand{\fracone}{1}
  \newcommand{\frachtwo}{5}
  \newcommand{\datacsvscatterone}{plots/scatter/scatter_results_cifar/ind_0_dropfrac_0.01.csv}
  \newcommand{\datacsvscattertwo}{plots/scatter/scatter_results_cifar/ind_0_dropfrac_0.05.csv}
  \pgfplotsset{colormap/Set1}

\definecolor{Set1Red}{RGB}{74,124,182}
\definecolor{Set1Blue}{RGB}{142,82,159}
\definecolor{Set1Green}{RGB}{239,134,50}

\usetikzlibrary{calc}

\def\datacsvcorrelations{plots/scatter/correlations_\dataset/dropfrac_\fracone/example_\exampleid.csv}

\def\dataset{wikitext}
\def\datacsvcorrelations{plots/scatter/correlations_\dataset/dropfrac_\frachtwo/example_\exampleid.csv}

\pgfplotstableread[col sep=comma]{\datacsvcorrelations}\datatable

\pgfplotstablegetelem{0}{corr}\of\datatable
\pgfmathprintnumberto[fixed, precision=2]{\pgfplotsretval}{\trakcorrhtwo}

\pgfplotstablegetelem{1}{corr}\of\datatable
\pgfmathprintnumberto[fixed, precision=2]{\pgfplotsretval}{\ekfaccorrhtwo}

\pgfplotstablegetelem{2}{corr}\of\datatable
\pgfmathprintnumberto[fixed, precision=2]{\pgfplotsretval}{\methodcorrhtwo}


\ifdefined\hidelegend
\pgfplotsset{conditional legend style/.style={legend
style={draw=none, fill=none}}}
\else
\pgfplotsset{conditional legend style/.style={
    legend style={
      at={(1.95, 1.45)}, %
      anchor=north east, %
      legend columns=3, %
      /tikz/every even column/.append style={column sep=0.5cm}, %
      inner xsep=5pt, %
      inner ysep=2.5pt,  %
      legend image post style={scale=2.5} %
    }
}}
\fi

\begin{tikzpicture}
  \begin{axis}[
      name=bar_chart, %
      ybar,           %
      width=0.27\textwidth, %
      height=0.3\textwidth, %
      ylabel={Spearman $\rho$},
      xlabel={Method},
      symbolic x coords={A, B, C}, %
      xtick=data,     %
      xticklabels={}, %
      xticklabel style={rotate=45, anchor=east}, %
      nodes near coords, %
      nodes near coords align={vertical},
      nodes near coords style={color=gray}, %
      enlarge x limits=1.5, %
      ymin=0,           %
      ymax=1.2,        %
      axis lines*=left,  %
      bar width=15pt,   %
      title={Correlations}, %
    ]
    \addplot +[fill=Set1Red, draw=none, opacity=0.5] coordinates {(A, 0.234)};
    \addplot +[fill=Set1Blue, draw=none, opacity=0.5] coordinates {(B, 0.347)};
    \addplot +[fill=Set1Green, draw=none] coordinates {(C, 0.957)};
  \end{axis}

  \begin{axis}[
      name=plot1,
      at=(bar_chart.east), %
      anchor=west,         %
      xshift=1.5cm,        %
      width=0.32\textwidth,
      height=0.3\textwidth,
      xlabel={Predicted Loss},
      ylabel={True Loss},
      title={Rand. CIFAR-10 subsets},
      conditional legend style,
      grid=both,
      grid style={line width=.1pt, draw=gray!10},
      major grid style={line width=.2pt,draw=gray!50},
      scaled y ticks=true,
      yticklabel style={/pgf/number format/.cd, fixed, precision=3},
      xticklabel style={/pgf/number format/.cd, fixed, precision=3},
      clip=false,
      after end axis/.code={
        \pgfmathsetmacro{\myxmin}{\pgfkeysvalueof{/pgfplots/xmin}}
        \pgfmathsetmacro{\myxmax}{\pgfkeysvalueof{/pgfplots/xmax}}
        \pgfmathsetmacro{\myymin}{\pgfkeysvalueof{/pgfplots/ymin}}
        \pgfmathsetmacro{\myymax}{\pgfkeysvalueof{/pgfplots/ymax}}
        \pgfmathsetmacro{\linestart}{max(\myxmin, \myymin)}
        \pgfmathsetmacro{\lineend}{min(\myxmax, \myymax)}
        \draw[dashed, gray, very thick, opacity=0.5]
        (axis cs:\linestart,\linestart) --
        (axis cs:\lineend,\lineend);
      }
    ]
    \addplot[
      only marks,
      mark=square*,
      mark size=1.5pt,
      color=Set1Red,
      fill opacity=\scatteropacity,
      draw opacity=\scatteropacity,
      point meta=explicit,
    ] table[x=kron_infl_10, y=true, meta expr=1, col sep=comma]
    {\datacsvscatterone};
    \ifdefined\hidelegend\else
    \addlegendentry{EK-FAC (re-scaled)}
    \fi

    \addplot[
      only marks,
      mark=triangle*,
      mark size=1.5pt,
      color=Set1Blue,
      fill opacity=\scatteropacity,
      draw opacity=\scatteropacity,
      point meta=explicit,
    ] table[x=trak_infl_10, y=true, meta expr=2, col sep=comma]
    {\datacsvscatterone};
    \ifdefined\hidelegend\else
    \addlegendentry{TRAK (re-scaled)}
    \fi

    \addplot[
      only marks,
      mark=*,
      mark size=2pt,
      color=Set1Green,
      fill opacity=1,
      draw opacity=1,
      point meta=explicit,
    ] table[x=infls, y=true, meta expr=0, col sep=comma] {\datacsvscatterone};
    \ifdefined\hidelegend\else
    \addlegendentry{\method{} (unscaled)}
    \fi

    \ifdefined\hidelegend\else
    \addlegendimage{dashed, gray, very thick,
      legend image code/.code={
        \draw[dashed, gray, very thick] (0cm,0cm) -- (0.32cm,0cm); %
      }
    }
    \fi

    \node (annotation_text) at (axis cs: 2.3, 1.9) [align=center,
    text width=4.5cm] {For each color, each dot is a training data subset};
    \coordinate (target_point) at (axis cs: 1.682, 1.87);
    \draw [->, thick, gray, shorten >=1pt] (annotation_text.west) --
    (target_point);

    \node (annotation_text) at (axis cs: 2.32, 1.6) [align=center,
    text width=6.2cm] {\method{} nearly {\em perfectly} predicts loss};
    \coordinate (target_point) at (axis cs: 1.625, 1.6);
    \draw [->, thick, Set1Green, shorten >=1pt]
    (annotation_text.west) -- (target_point);
    \coordinate (target_point_2) at (axis cs: 1.55, 1.54);
    \draw [->, thick, Set1Green, shorten >=1pt]
    (annotation_text.west) -- (target_point_2);
    \coordinate (target_point_3) at (axis cs: 1.827, 1.83);
    \draw [->, thick, Set1Green, shorten >=1pt]
    (annotation_text.west) -- (target_point_3);

    \node (annotation_text) at (axis cs: 2.3, 1.74) [align=center,
    text width=6.2cm] {State-of-the-art attribution methods\\  only
    {\em weakly} correlate with loss};
    \coordinate (target_point) at (axis cs: 1.74, 1.62);
    \draw [->, thick, Set1Red, opacity=0.5, shorten >=1pt]
    (annotation_text.west) -- (target_point);
    \coordinate (target_point_1) at (axis cs: 1.8, 1.75);
    \draw [->, thick, Set1Blue, opacity=0.5, shorten >=1pt]
    (annotation_text.west) -- (target_point_1);
    \coordinate (target_point_3) at (axis cs: 1.683, 1.83);
    \draw [->, thick, Set1Red, opacity=0.5, shorten >=1pt]
    (annotation_text.west) -- (target_point_3);
  \end{axis}
\end{tikzpicture}

  \caption{ \method{} nearly perfectly predicts the effect of training data
    removal. In contrast to the best baselines \citep{park2023trak,grosse2023studying}, 
    \method{} produces estimates
    that both (a) highly correlate with the ground truth effect and (b) are
    well-scaled. {\bf Right:} we plot the predicted loss (from \method{} and 
    the two baselines) against the true loss for a
    randomly chosen test point, each point a training data subset with a random
    1\% of samples removed. For \method{} we plot the predicted loss directly 
    since it is well-scaled; for the baselines, we first rescale the predictions 
    to match the variance of the ground-truth losses.
    {\bf Left:} The average (taken across test examples) 
    Spearman correlation between predicted and true model losses (also known as 
    the LDS \citep{ilyas2022datamodels,park2023trak}, see Section~\ref{sec:single_model_data_attribution}).}
  \label{fig:headline}
\end{figure}

\paragraph{A new perspective (single-model data
attribution).} The inherent randomness of large-scale training makes
attributing specific model predictions to training data
conceptually
challenging \citep{bae2022if,ilyas2022datamodels,nguyen2023bayesian}. After all,
if training on the same data can lead to different models, then
we {\em cannot} predict the variation between these models as a
function of the dataset.
As a result, prior work can only predict how a {\em learning algorithm}
would behave ({\em on average}) if trained on a different dataset,
but not how a {\em specific} model would behave if the training data were
different.

Motivated by this state of affairs,
we introduce a setting called ``single-model'' data attribution.
The goal in this setting is still to predict the behavior of a model under
changes to the training data---the twist is that
we aim to predict how {\em this specific model} would have behaved
under different training data, rather than how a newly initialized
and trained model would have behaved
under different training data.
This subtle change means that: (a) in the single-model setting,
it {\em is} possible to perfectly predict model behavior as a function
of the training data, and (b)
these predictions correspond to how a given ``single model''
would respond to changes in the training data (motivating the name
of the setting), rather than a given learning algorithm.



\paragraph{A new data attribution method (\method{}).}
We present \method{}
(\textbf{M}etagradient-based
  \textbf{A}ttribution via
  \textbf{G}round-truth
  \textbf{I}nfluence
\textbf{C}omputation),
a state-of-the-art
data attribution method.
Our method leverages recent advances in large-scale metagradient calculation
\citep{engstrom2025optimizing} to \textit{exactly} calculate the
influence function \citep{hampel1947influence} for large-scale learning
algorithms.
\method{} accurately estimates how
model predictions respond to random training data deletions---substantially
outperforming existing methods---\textit{even in our more challenging
single-model setting}.
For example (see Figure~\ref{fig:headline}),
\begin{itemize}
  \item When dropping different random 1\% subsets from the training set of a
    ResNet-9 model trained on CIFAR-10, \method{} almost exactly predicts
    ground-truth model losses ($\rho = 0.96$) while existing
    methods' predictions are only weakly correlated ($\rho = 0.25$).
  \item When dropping different random 1\% subsets
    from the training set of a Gemma-2B model trained on instruction tuning
    data, \method{} nearly exactly predicts ground-truth model test losses
    ($\rho = 0.97$) while existing methods perform no better
    than random guessing.
\end{itemize}

\noindent{}Together, the single-model data attribution setting and our new
primitive, \method{}, enable nearly optimally estimating the effect of removing
and adding training data on trained model predictions in modern (deep learning) 
settings, including language modeling and supervised vision tasks.






